
\subsubsection{3.1 Problem Formulation}

Consider distilling a pretrained Transformer teacher T into an SSM-based student S. At layer $\ell$, the teacher computes attention output via softmax attention over query, key, and value projections. The student computes SSM output via selective state space dynamics.

\textbf{Matrix-level (MOHAWK)}: Minimize attention map discrepancy:

$$\mathcal{L}_{\text{matrix}} = \sum_{\ell} \left\| \text{Mixer}_T^{(\ell)}(\mathbf{x}) - \text{Mixer}_S^{(\ell)}(\mathbf{x}) \right\|_F^2$$

\textbf{Token-level (CAB)}: Align individual projections via learned bridges:

$$\mathcal{L}_{\text{token}} = \sum_{\ell} \left\| f_B(\mathbf{Q}^{(\ell)}) - \mathbf{B}_S^{(\ell)} \right\|^2 + \left\| f_C(\mathbf{K}^{(\ell)}) - \mathbf{C}_S^{(\ell)} \right\|^2$$

where $f_B, f_C$ are learned MLP bridges mapping Transformer projections to SSM parameters.

The key distinction: matrix-level objectives supervise position-position relationships (the attention map); token-level objectives supervise individual token representations.

\subsubsection{3.2 Unified Framework Design}

Both objectives were implemented in a single Phi-Mamba codebase ensuring:

\begin{itemize}
\item Same teacher model: Phi-1.5 (1.3B parameters)
\item Same student architecture: Phi-Mamba with MOHAWK modifications (multi-head SSM, no $\Delta$, open gates)
\item Same training data: C4 dataset (streaming), Phi tokenizer
\item Same optimization: AdamW, identical learning rate schedules
\end{itemize}

Framework validation (H-E1) confirmed both objectives execute without errors in this unified setup.

\subsubsection{3.3 Hidden State Drift Measurement}

To quantify representation stability across lengths, hidden state drift at layer $\ell$ for sequence length L is defined as:

$$\text{Drift}^{(\ell)}(L) = \frac{1}{N} \sum_{i=1}^{N} \left\| \mathbf{h}_T^{(\ell)}(x_i, L) - \mathbf{h}_S^{(\ell)}(x_i, L) \right\|_2$$

where $\mathbf{h}_T, \mathbf{h}_S$ are teacher and student hidden states, and the average is over N samples.

\textbf{Drift slope}: Linear regression of Drift(L) against L yields slope indicating how quickly representations diverge with length. Lower slope indicates more stable representations.

\textbf{Drift ratio}: Drift(L\_max) / Drift(L\_min) measures relative degradation. Ratio < 2.0 indicates bounded drift.

